{{دیوان
| عنوان = ہمدردی
| ماخذ = ویکی ماخذ
| ماخذ_لنک = https://ur.wikisource.org/wiki/%DB%81%D9%85%D8%AF%D8%B1%D8%AF%DB%8C
| مقابلہ = اقبال سائبر لائبریری (اقبال اکادمی پاکستان)
| مقابلہ_لنک = https://www.iqbalcyberlibrary.net
}}
=== (ماخوذ از ولیم کوپر) ===

=== بچوں کے لیے ===

<poem>
ٹہنی پہ کسی شجر کی تنہا
بُلبل تھا کوئی اُداس بیٹھا

کہتا تھا کہ رات سر پہ آئی
اُڑنے چُگنے میں دن گزارا

پہنچوں کِس طرح آشیاں تک
ہر چیز پہ چھا گیا اندھیرا

سُن کر بُلبل کی آہ و زاری
جُگنو کوئی پاس ہی سے بولا

حاضر ہُوں مدد کو جان و دل سے
کِیڑا ہوں اگرچہ مَیں ذرا سا

کیا غم ہے جو رات ہے اندھیری
مَیں راہ میں روشنی کروں گا

اللہ نے دی ہے مجھ کو مشعل
چمکا کے مجھے دِیا بنایا

ہیں لوگ وہی جہاں میں اچھّے
آتے ہیں جو کام دوسرں کے
</poem>
